\documentclass[11pt]{article}
\usepackage[a4paper,margin=1in]{geometry}
\usepackage{amsmath,amssymb,amsthm}
\usepackage{array}
\usepackage{hyperref}

\newcommand{\dd}{\mathrm{d}}
\newcommand{\sgn}{\operatorname{sgn}}
\newcommand{\He}{\mathrm{He}}
\newcommand{\cA}{\mathcal{A}}
\newcommand{\avg}[1]{\langle #1\rangle}
\newcommand{\vt}{v_t}
\renewcommand{\Re}{\operatorname{Re}}
\renewcommand{\Im}{\operatorname{Im}}

\newtheorem{theorem}{Theorem}
\newtheorem{lemma}[theorem]{Lemma}
\newtheorem{corollary}[theorem]{Corollary}
\theoremstyle{remark}
\newtheorem{remark}[theorem]{Remark}

\title{The exact linear heat-flux closure for the conserving\\
Lenard--Bernstein (Dougherty) operator, in one and three velocity dimensions}
\author{Working note}
\date{16 November 2026}

\begin{document}
\maketitle

\noindent\emph{Summary.} The argument of
\texttt{bgk\_velocity\_dependent\_closure.tex} is carried over to the
conserving Lenard--Bernstein (Dougherty) operator, which is a Fokker--Planck
operator and is the ground truth of the tournament note. Three things are
proved.
\begin{enumerate}
  \item The dual variational principles hold for any conserving, symmetric,
    non-negative linearized collision operator in one velocity dimension
    (Theorem~\ref{thm:abs}).
  \item For the Dougherty operator in one velocity dimension the
    quasi-static closure has a closed form, with sharp bounds
    (Theorem~\ref{thm:1d}).
  \item For the Dougherty operator in three velocity dimensions, which is
    the tournament's operator, the closure also has a closed form
    (Theorem~\ref{thm:3d}). It is positive, strictly below the
    Chapman--Enskog conductivity $\tfrac56\,n_0T_0/(m\nu)$, departs from it
    as $\tfrac{2}{15}(k\lambda)^2$, and tends to the collisionless value.
\end{enumerate}
The Landau (Coulomb) operator is not treated; Section~\ref{sec:landau} says
what stands in the way. A consequence for the benchmark is stated in
Section~\ref{sec:impl}: with a velocity-independent collision frequency the
local closure does not fail until $k\lambda$ is of order one. No second
reader has checked the proofs.

\section{Abstract variational theorem (one velocity dimension)}

Units and notation as in the previous notes: $M$ the unit Maxwellian,
$\avg{\cdot}$ its average, $H = \operatorname{span}\{1,v,v^2\}$.

\paragraph{Assumption (A).} $\Lambda$ is a self-adjoint, non-negative
operator on $L^2(M\,\dd v)$ that is real, preserves parity, and has kernel
$H$; and for $k\ne0$ the operator $\cA_k = ikv+\Lambda$ is invertible, with
$(-i\omega+\cA_k)^{-1}$ continuous in $\omega$ near $0$ in $\Im\omega\ge0$.

For a real function $\chi$ put
$\avg{\chi,\Lambda^+\chi} = \sup_f\bigl[2\avg{\chi f}-\avg{f\Lambda f}\bigr]
\in[0,\infty]$, the supremum over real $f$ in the form domain of $\Lambda$.
Then $\avg{\chi f}^2\le\avg{\chi,\Lambda^+\chi}\,\avg{f\Lambda f}$, and if
$\Lambda f_0=\chi$ then $\avg{\chi,\Lambda^+\chi} = \avg{f_0\Lambda f_0}$.

\begin{theorem}\label{thm:abs}
Assume (A), and let $g = \cA_k^{-1}1 = g_e+iw$.
\begin{enumerate}
  \item[(i)] $\avg{v^2g}=0$, $n_1=\avg{g}>0$, and the quasi-static limit of
    the potential-driven response is
    $\hat q = -i\,\sgn(k)\,\theta\,\hat T$ with $\theta = 2/(|k|n_1)$.
  \item[(ii)] For even $\varphi$ with $\avg{v^2\varphi}=0$ and odd $\psi$
    with $k\avg{v\psi}=k\avg{v^3\psi}=1$,
    \[
      \frac{\avg{\varphi}^2}
           {\avg{\varphi\Lambda\varphi}+k^2\avg{v\varphi,\Lambda^+v\varphi}}
      \;\le\; n_1 \;\le\;
      \avg{\psi\Lambda\psi}+\avg{(1-kv\psi),\Lambda^+(1-kv\psi)},
    \]
    with equality at $\varphi=g_e$, $\psi=-w$.
  \item[(iii)] With $\xi$ a solution of $\Lambda\xi=\He_3$ normalized by
    $\avg{v\xi}=0$,
    $\kappa_{\mathrm F} = \tfrac12\avg{\He_3\,\xi}$,
    $\chi_* = 1-v^2+v\xi/\kappa_{\mathrm F}$ and
    $K_2 = \tfrac12\kappa_{\mathrm F}\avg{\chi_*,\Lambda^+\chi_*}$,
    \[
      \frac{\kappa_{\mathrm F}}{1+K_2k^2}\;\le\;\frac{\theta}{|k|}
      \;\le\;\kappa_{\mathrm F}.
    \]
\end{enumerate}
\end{theorem}

\begin{proof}
(i) and the equality cases of (ii) are proved exactly as in the previous
note, which used only (A): pairing $\cA_kg=1$ with $1,v,v^2$; the real system
$-kvw+\Lambda g_e=1$, $kvg_e+\Lambda w=0$; and
$n_1=\avg{g_e\Lambda g_e}+\avg{w\Lambda w}$, where now
$\avg{w\Lambda w} = k^2\avg{vg_e,\Lambda^+vg_e}$ and
$\avg{g_e\Lambda g_e} = \avg{(1+kvw),\Lambda^+(1+kvw)}$.

Weak duality: with $\chi=1-kv\psi$,
\begin{align*}
  \avg{\varphi} &= \avg{\varphi\chi}+k\avg{v\varphi\,\psi} \\
  &\le\avg{\chi,\Lambda^+\chi}^{1/2}\avg{\varphi\Lambda\varphi}^{1/2}
  + \bigl(k^2\avg{v\varphi,\Lambda^+v\varphi}\bigr)^{1/2}
    \avg{\psi\Lambda\psi}^{1/2},
\end{align*}
and the Cauchy--Schwarz inequality in $\mathbb{R}^2$ finishes it.

(iii) Upper bound: $\varphi_0=(3-v^2)/2$ has $\Lambda\varphi_0=0$,
$v\varphi_0=-\He_3/2$, so the left side of (ii) is
$1/(\tfrac14k^2\avg{\He_3\xi}) = 2/(k^2\kappa_{\mathrm F})$. Lower bound:
$\psi_*=(v-\xi/\kappa_{\mathrm F})/k$ satisfies $k\avg{v\psi_*}=1$ and
$k\avg{v^3\psi_*} = 3-\avg{\He_3\xi}/\kappa_{\mathrm F}=1$, with
$\avg{\psi_*\Lambda\psi_*} = \avg{\xi\He_3}/(k\kappa_{\mathrm F})^2
= 2/(k^2\kappa_{\mathrm F})$ and $1-kv\psi_*=\chi_*$.
\end{proof}

\section{The Dougherty operator in one velocity dimension}

Let $N=-\partial_v^2+v\,\partial_v$, so $N\He_n=n\He_n$. The linearized
conserving Lenard--Bernstein operator is
\begin{equation}
  \Lambda h = \nu\bigl(Nh-\avg{vh}\,v-\avg{\He_2h}\,\He_2\bigr),
\end{equation}
that is, $\Lambda\He_n=\nu n\He_n$ for $n\ge3$ and $\Lambda\He_n=0$ for
$n\le2$. It satisfies (A): $\cA_k$ is a finite-rank perturbation of
$\cA^0_k = ikv+\nu N$, which has compact resolvent and trivial kernel, and
$\cA_k$ is injective by the argument of the previous note.

For $a>0$ let
\begin{equation}
  E_a(s) = \bigl[(1-s)\,e^{s}\bigr]^{a},\qquad 0\le s\le1 .
\end{equation}
$E_a$ decreases from $1$ to $0$ and $E_a(s)\le e^{-as^2/2}$.

\begin{theorem}\label{thm:1d}
For the Dougherty operator in one velocity dimension, with $y=\nu/|k|$ and
$a=1/y^2$,
\begin{equation}
  \theta_1(y) = \frac{2}{y}\;
  \frac{\int_0^1 s\,E_a(s)\,\dd s}{\int_0^1E_a(s)\,\dd s} .
  \label{eq:th1}
\end{equation}
Moreover, with $\kappa_{\mathrm F}=1/\nu$, $\lambda=1/\nu$ and
$\chi=\theta_1/|k|$:
\begin{enumerate}
  \item[(a)] $0<\chi(k)<\kappa_{\mathrm F}$ for all $k\ne0$;
  \item[(b)] $\chi(k)\ge\kappa_{\mathrm F}/(1+\tfrac13(k\lambda)^2)$, and the
    constant $\tfrac13$ is sharp:
    $1-\chi/\kappa_{\mathrm F} = \tfrac13(k\lambda)^2+O((k\lambda)^4)$;
  \item[(c)] $\theta_1(y)\to\sqrt{8/\pi}$ as $y\to0$.
\end{enumerate}
\end{theorem}

\begin{lemma}\label{lem:semi}
Let $\cA^0_k=ikv+\nu N$ and $\sigma\ge0$. Then
$f_\sigma=(\cA^0_k+\sigma\nu)^{-1}1$ satisfies
\[
  \nu\avg{f_\sigma} = \int_0^1(1-s)^{a-1+\sigma}e^{as}\,\dd s,
  \qquad a=k^2/\nu^2 .
\]
For $\sigma=0$: $\nu\avg{f_0}=\tfrac1a+\int_0^1E_a$,
$\avg{vf_0}=1/(ik)$ and $\avg{v^2f_0}=\nu/k^2$.
\end{lemma}

\begin{proof}
Let $c=ik/\nu$, $A(t)=-c\,(1-e^{-\nu t})$ and $B(t)=\nu\int_0^tA^2$. Then
$h(t,v)=e^{A(t)v+B(t)}$ solves $\partial_th+\cA^0_kh=0$, $h(0)=1$, since
$A'+ik+\nu A=0$ and $B'=\nu A^2$. $A$ is imaginary, so $|h|=e^{B}$, and
$B=c^2[\nu t-2(1-e^{-\nu t})+\tfrac12(1-e^{-2\nu t})]$ tends to $-\infty$
linearly; hence $f_\sigma=\int_0^\infty e^{-\sigma\nu t}h\,\dd t$ converges
in $L^2(M)$ and solves $(\cA^0_k+\sigma\nu)f_\sigma=1$. Also
$\avg{h}=e^{B+A^2/2}=\exp\bigl(-a(\nu t-1+e^{-\nu t})\bigr)$. Substituting
$s=1-e^{-\nu t}$ gives the formula. For $\sigma=0$, integrate by parts:
$\int_0^1(1-s)^{a-1}e^{as}\dd s=\tfrac1a+\int_0^1(1-s)^ae^{as}\dd s$. The
last two identities follow by pairing $\cA^0_kf_0=1$ with $1$ and $v$:
$ik\avg{vf_0}=1$ and $ik\avg{v^2f_0}+\nu\avg{vf_0}=0$.
\end{proof}

\begin{proof}[Proof of Theorem~\ref{thm:1d}]
Write $S=\int_0^1E_a$. By Theorem~\ref{thm:abs}(i),
$\avg{vg}=1/(ik)$ and $\avg{\He_2g}=-n_1$, so $\cA_kg=1$ reads
\[
  \cA^0_kg = 1+\frac{\nu}{ik}\,v-\nu n_1\He_2 .
\]
$\cA^0_k$ is symmetric for the bilinear pairing $\avg{fh}$, so
$\avg{(\cA^0_k)^{-1}\varrho} = \avg{\varrho f_0}$. By Lemma~\ref{lem:semi},
$\avg{vf_0}=1/(ik)$ and $\avg{\He_2f_0}=\nu/k^2-\avg{f_0}$. Hence
\[
  n_1 = \avg{f_0}-\frac{\nu}{k^2}-\nu n_1\Bigl(\frac{\nu}{k^2}-\avg{f_0}\Bigr),
  \qquad\text{so}\qquad \nu n_1 = \frac{S}{1-S},
\]
using $\nu\avg{f_0}-\nu^2/k^2=S$. Thus $\theta_1=2y(1-S)/S$. Finally
$E_a' = -a\,s\,E_a/(1-s)$, and integrating by parts,
$1-S=\int_0^1(1-E_a) = -\int_0^1(1-s)E_a'\,\dd s = a\int_0^1sE_a$, which
gives \eqref{eq:th1}.

(a) Positivity is evident. Since $\int_0^1(s-\tfrac12)\dd s=0$ and $E_a$ is
strictly decreasing,
$\int_0^1(s-\tfrac12)E_a = \int_0^1(s-\tfrac12)(E_a(s)-E_a(\tfrac12))\dd s<0$,
so $\theta_1<1/y$.

(b) Apply Theorem~\ref{thm:abs}(iii): $\xi=\He_3/(3\nu)$,
$\kappa_{\mathrm F}=\tfrac12\avg{\He_3^2}/(3\nu)=1/\nu$,
$\chi_*=1-v^2+v\He_3/3=\He_4/3$, and
$\avg{\chi_*,\Lambda^+\chi_*}=\avg{\He_4^2}/(9\cdot4\nu)=2/(3\nu)$, so
$K_2=1/(3\nu^2)$. Sharpness: with $L(s)=s+\ln(1-s)\le0$,
$\int_0^1L=-\tfrac12$ and $\int_0^1L^2=\tfrac56$, so
$S=1-\tfrac a2+\tfrac{5}{12}a^2+O(a^3)$ and
$y\theta_1=2(1-S)/(aS)=1-\tfrac a3+O(a^2)$.

(c) Substituting $s=u/\sqrt a$ and using dominated convergence with the
bound $E_a\le e^{-as^2/2}$: $\sqrt a\int_0^1E_a\to\sqrt{\pi/2}$ and
$a\int_0^1sE_a\to1$. Hence
$\theta_1\to(2/y)\,a^{-1}\sqrt{2a/\pi}=\sqrt{8/\pi}$.
\end{proof}

\section{The Dougherty operator in three velocity dimensions}

This is the operator of the tournament note, linearized about a uniform
Maxwellian. With $v=(v_x,v_y,v_z)$, $M$ the unit Maxwellian on
$\mathbb{R}^3$, $N=-\Delta_v+v\cdot\nabla_v$, and
$\mathcal{E}=(|v|^2-3)/\sqrt6$,
\begin{equation}
  \Lambda h = \nu\Bigl(Nh-\sum_{j}\avg{v_jh}\,v_j
     -2\avg{\mathcal{E}h}\,\mathcal{E}\Bigr),
\end{equation}
with spatial dependence on $x$ only and force along $x$. The moments are
\[
  T=\tfrac13\avg{(|v|^2-3)h},\qquad
  q=\tfrac12\avg{v_x(|v|^2-5)h},
\]
and the Chapman--Enskog conductivity is
$\kappa_{\mathrm F}=\tfrac56\,n_0T_0/(m\nu)$.

\begin{theorem}\label{thm:3d}
For the three-dimensional Dougherty operator the quasi-static limit of the
potential-driven response is $\hat q=-i\,\sgn(k)\,\theta_3(\nu/|k|)\,\hat T$
with
\begin{equation}
  \theta_3(y) = \frac1y\;
  \frac{\int_0^1s\,(3-2s)\,E_a(s)\,\dd s}{\int_0^1E_a(s)\,\dd s},
  \qquad a=1/y^2 .
  \label{eq:th3}
\end{equation}
With $\lambda=\vt/\nu$ (here $\vt=1$) and $\chi=\theta_3/|k|$:
\begin{enumerate}
  \item[(a)] $0<\chi(k)<\kappa_{\mathrm F}$ for all $k\ne0$;
  \item[(b)] $1-\chi/\kappa_{\mathrm F}
     =\tfrac{2}{15}(k\lambda)^2+O((k\lambda)^4)$ as $k\lambda\to0$;
  \item[(c)] $\theta_3(y)\to3\sqrt{2/\pi}$ as $y\to0$, which is
    $\tfrac32$ times the one-dimensional Hammett--Perkins value, as it must
    be when only the parallel temperature responds and $T=T_\parallel/3$.
\end{enumerate}
\end{theorem}

\begin{proof}
\emph{Reduction.} Let $\xi=(v_\perp^2-2)/2$ with $v_\perp^2=v_y^2+v_z^2$;
$\avg{\xi}=0$, $\avg{\xi^2}=1$, $N\xi=2\xi$, and
$|v|^2-3=\He_2(v_x)+2\xi$. Streaming acts on $v_x$ only, so $\cA_k$
preserves the space of functions $c(v_x)+d(v_x)\,\xi$, on which
$\cA^0_k=ikv_x+\nu N$ acts as $\cA^x_k$ on $c$ and as $\cA^x_k+2\nu$ on $d$,
where $\cA^x_k=ikv_x+\nu N_x$ is the one-dimensional operator of
Lemma~\ref{lem:semi}. Invertibility of $\cA_k$ on this space follows as in
one dimension.

\emph{Moment relations.} Let $g=\cA_k^{-1}1$ and $n_1=\avg{g}$. Pairing
$\cA_kg=1$ with $1$, $v_x$ and $|v|^2$:
$ik\avg{v_xg}=1$, $ik\avg{v_x^2g}=0$, $ik\avg{v_x|v|^2g}=3$. Hence
$q_1=\tfrac12\avg{v_x(|v|^2-5)g}=(3-5)/(2ik)=i/k$,
$\avg{\He_2(v_x)g}=-n_1$, and with $\epsilon_1=\avg{\xi g}$,
\begin{equation}
  3T_1 = -n_1+2\epsilon_1 .
  \label{eq:T1}
\end{equation}

\emph{Solution.} Since
$2\avg{\mathcal{E}g}\mathcal{E}=T_1(|v|^2-3)=T_1(\He_2(v_x)+2\xi)$,
\[
  \cA^0_kg = 1+\frac{\nu}{ik}\,v_x+\nu T_1\He_2(v_x)+2\nu T_1\,\xi .
\]
With $f_0$, $f_2$ as in Lemma~\ref{lem:semi} (for $\sigma=0,2$) and
$S_j=\int_0^1(1-s)^jE_a$: $\nu\avg{f_0}-\nu^2/k^2=S_0$ and
$\nu\avg{f_2}=S_1$. Taking the average, and the average against $\xi$,
\[
  \nu n_1 = S_0\,(1-\nu T_1),\qquad \epsilon_1 = 2\,T_1\,S_1 .
\]
Substituting into \eqref{eq:T1} with $\tau_1=\nu T_1$:
$\tau_1\,(3-S_0-4S_1)=-S_0$. Therefore, for $k>0$,
\[
  \frac{q_1}{T_1} = \frac{i\nu}{k\,\tau_1}
  = -\,i\,y\,\frac{3-S_0-4S_1}{S_0}.
\]
Finally $3-S_0-4S_1=\int_0^1(5-4s)(1-E_a)\,\dd s$, and integrating by parts
with the antiderivative $-(1-s)(3-2s)$ of $5-4s$ gives
$a\int_0^1s(3-2s)E_a$. This is \eqref{eq:th3}; the case $k<0$ follows by
symmetry.

(a) Positivity is evident. Let $p(s)=s(3-2s)$; $\int_0^1p=\tfrac56$, and
$p-\tfrac56$ is negative on $[0,s_0)$ and positive on $(s_0,1]$, with
$s_0=(9-\sqrt{21})/12$. Since $E_a$ is strictly decreasing,
$\int_0^1(p-\tfrac56)E_a=\int_0^1(p-\tfrac56)(E_a-E_a(s_0))<0$. Hence
$\theta_3<5/(6y)$, that is, $\chi<\kappa_{\mathrm F}$.

(b) For $x\le0$, $|e^x-1-x|\le x^2/2$. With $L=s+\ln(1-s)$:
$\int_0^1L=-\tfrac12$ and $\int_0^1pL=-\tfrac{19}{36}$, so
$\int E_a=1-\tfrac a2+O(a^2)$,
$\int pE_a=\tfrac56-\tfrac{19}{36}a+O(a^2)$, and
$\tfrac65y\theta_3=1-\tfrac{2}{15}a+O(a^2)$.

(c) As in Theorem~\ref{thm:1d}(c), $\sqrt a\int E_a\to\sqrt{\pi/2}$ and
$a\int pE_a\to3$.
\end{proof}

\section{Numerical checks}

Scratch runs on 16 November 2026, not part of the repository.
\begin{itemize}
  \item \textbf{Independent solves.} $\cA_kg=1$ was solved in a Hermite
    basis (one channel in one dimension; two channels, for $1$ and $\xi$, in
    three). For $y$ from $0.02$ to $30$ the result agrees with
    \eqref{eq:th1} and \eqref{eq:th3} to nine digits.
  \item \textbf{Bounds.} $0<\theta_1<1/y$ and $0<\theta_3<5/(6y)$ on grids
    over several decades. $y^2(1-y\theta_1)$ rises to $0.3329$ at $y=20$ and
    $y^2(1-\tfrac65y\theta_3)$ to $0.1332$, against $\tfrac13$ and
    $\tfrac{2}{15}$.
  \item \textbf{A numerical trap.} Evaluating $\theta_1$ from
    $2y(1+y^2-\mu)/(\mu-y^2)$ with
    $\mu=\int_0^1(1-s)^{a-1}e^{as}\dd s$, which is algebraically the same,
    loses all accuracy for $y\gtrsim100$ and appeared to violate (b). The
    forms \eqref{eq:th1} and \eqref{eq:th3} do not have this problem.
\end{itemize}
\begin{center}
\renewcommand{\arraystretch}{1.2}
\begin{tabular}{rrrr}
  $k\lambda$ & $\chi/\kappa_{\mathrm F}$ (3D) & $1-\tfrac{2}{15}(k\lambda)^2$
     & $\chi/\kappa_{\mathrm F}$ (1D) \\ \hline
  0.1 & 0.9987 & 0.9987 & 0.9967 \\
  0.3 & 0.9883 & 0.9880 & 0.9715 \\
  1   & 0.8946 & 0.8667 & 0.7844 \\
  3   & 0.5940 & --- & 0.4165 \\
  10  & 0.2467 & --- & 0.1484 \\
\end{tabular}
\end{center}

\paragraph{Observed but not proved.} $\theta_1$ and $\theta_3$ are
decreasing and stay below their collisionless values. In three dimensions
$1-\chi/\kappa_{\mathrm F}\le\tfrac{2}{15}(k\lambda)^2$ on the grid; only the
asymptotic statement (b) is proved. The harmonic-mean interpolant
$(1/\theta(0)+|k|^{-1}\kappa_{\mathrm F}^{-1})^{-1}$ lies below the exact
value by at most 22\% in one dimension and 19\% in three.

\section{Consequences for the benchmark}
\label{sec:impl}

\begin{itemize}
  \item \textbf{An exact reference for the tournament's ground truth.} For
    linear, potential-driven, quasi-static tests, \eqref{eq:th3} is the
    answer at every collisionality. It gives the three checks any submitted
    closure should pass in that regime: the Chapman--Enskog limit with
    deficit $\tfrac{2}{15}(k\lambda)^2$, the collisionless limit
    $3\sqrt{2/\pi}$, and $\chi<\kappa_{\mathrm F}$ at every $k$.
  \item \textbf{The local closure fails late in this model.} With $\nu$
    independent of $v$, Fourier's law is within 1.2\% at $k\lambda=0.3$ and
    within 11\% at $k\lambda=1$. The early failure of Spitzer--H\"arm
    transport, at $k\lambda\sim10^{-2}$, is not present. The previous note
    showed that a $|v|^{-3}$ tail in the collision frequency moves the
    breakdown there, and reverses the sign of the error of harmonic-mean
    closures. The tournament note motivates the benchmark by that early
    failure; the Dougherty ground truth does not reproduce it. A
    velocity-dependent collision frequency would.
  \item \textbf{Sign of the interpolation error.} In all three
    velocity-independent models (BGK, Dougherty in one and three dimensions)
    the harmonic mean is too low; with the Coulomb-like tail it is too high,
    as reported for the Landau operator. A closure tuned on Dougherty data
    may therefore carry the wrong sign of error to a Coulomb plasma.
\end{itemize}

\section{What is not done: the Landau operator}
\label{sec:landau}

\begin{enumerate}
  \item \textbf{Closed form.} The Dougherty result rests on the explicit
    semigroup of Lemma~\ref{lem:semi}, which exists because the operator is
    diagonal in Hermite functions. The linearized Landau operator is not,
    and no closed form should be expected.
  \item \textbf{Variational bounds in three dimensions.}
    Theorem~\ref{thm:abs} is for one velocity dimension, where momentum
    conservation gives $\avg{v^2g}=0$ and turns $T_1$ into the diagonal
    resolvent element $-\avg{1,\cA_k^{-1}1}$. In three dimensions momentum
    conservation gives $\avg{v_x^2g}=0$, which involves the parallel viscous
    stress, and $T_1=\tfrac13\avg{(|v|^2-3)g}$ is an off-diagonal element.
    Positivity and the dual bounds then do not follow from the argument
    here. For the Dougherty operator they were recovered from the explicit
    formula. For Landau one would need a variational principle for the
    coupled heat-flux and stress response, which I have not found.
  \item \textbf{Inverting the operator on trial functions.} The bounds need
    $\Lambda^+$ applied to explicit functions. For BGK this is division by
    $\nu$ and for Dougherty division by the Hermite index. For Landau it is
    the Spitzer problem, solvable only numerically, although upper bounds on
    $\avg{\chi,\Lambda^+\chi}$ can be had from trial fluxes.
  \item \textbf{What is proved for a Coulomb-like collision frequency} is
    the relaxation-model result of the previous note, in one velocity
    dimension.
\end{enumerate}
Other limits are as before: linear response, a single drive, quasi-static,
unbounded or periodic domain. No prior-art search was made for the Dougherty
closed forms; exact solutions of Lenard--Bernstein models are classical and
these formulas may be known.

\end{document}
